%mac
\loai{Tính tốc độ trung bình}
\begin{vidu}
\immini[thm]{Hãy tính tốc độ trung bình ra $m/s$ và $km/h$ của nữ vận động viên tại một số giải thi đấu dựa vào số liệu trong bảng bên.}{\begin{tabular}{|c|c|c|}
\hline 
Giải thi đấu & Cự li chạy (m) & Thời gian chạy (s)\\ 
\hline 
Điền kinh quốc gia 2016&100  &11,64  \\ 
\hline 
SEA Games 29 (2017)&100  &11,56  \\ 
\hline 
SEA games 30 (2019)&100  &11,54  \\ 
\hline 
\end{tabular}} 
\haidong{12}
\loigiai{
Tốc độ trung bình của nữ vận động viên tại các giải thi đấu là:
\begin{itemize}
\item Điền kinh quốc gia 2016: $v_{tb}=\dfrac{100}{11,64}=8,59\ m/s=30,92\ km/h$.
\item SEA Games 29 (2017): $v_{tb}=\dfrac{100}{11,56}=8,65\ m/s=31,14\ km/h$
\item SEA Games 30 (2019): $v_{tb}=\dfrac{100}{11,54}=8,67\ m/s=31,21\ km/h$
\end{itemize}
}
\end{vidu}
\begin{vidu}
\immini[thm]{Bạn Mapi đi học từ nhà đến trường theo lộ trình $ABC$ như hình vẽ. Biết bạn Mapi đi đoạn đường $AB=200\ m$ mất $4$ phút, đoạn đường $BC=300\ m$ hết $5$ phút. Xác định tốc độ trung bình và vận tốc trung bình của bạn Mapi khi đi từ nhà đến trường.}
{\begin{tikzpicture}
\node[] at (0,0){\includegraphics[width=4.5cm]{images/5-2}};
\end{tikzpicture}}
\haidong{12}
\loigiai{
\begin{itemize}
\item Tốc độ trung bình của bạn Bình
\begin{itemize}
\item Quãng đường bạn Bình đi được: $s=AB+BC=500\ m$
\item Thời gian đi của bạn Bình là: $t=4+5=9$ phút $=50\ s$.
\item Tốc độ trung bình của bạn Bình là: $v_{tb}=\dfrac{s}{t}\approx 0,93\ m/s$.
\end{itemize}	
\item Vận tốc trung bình của bạn Bình:
\begin{itemize}
\item Chọn trục Ox trùng với AC, chiều dương của trục theo chiều từ A đến C.
\item Độ dịch chuyển là: $d=AC=\sqrt{AB^2+BC^2}\approx 360,6\ m$.
\item Vận tốc trung bình của bạn Bình là: $v=\dfrac{d}{t}\approx 0,67\ m/s$.
\end{itemize}	
\end{itemize}	
}
\end{vidu}
\begin{vidu}
Một con kiến bò quanh miệng của một cái chén được $1$ vòng hết $3$ giây. Bán kính của miệng chén là $3\ cm$.
\begin{itemize}
\item[a)] Tính quãng đường đi được và độ dịch chuyển của kiến. 
\item[b)] Tính tốc độ trung bình và vận tốc trung bình của con kiến ra $cm/s$.
\end{itemize}
\haidong{12}
\loigiai{
\begin{itemize}
\item[a)] Quãng đường đi được của con kiến là: $s=2\pi r=2\pi.0,03=0,06\pi\approx 0,1884\ m=18,84\ cm$.\\
Con kiến bò quanh miệng cốc, đi 1 vòng về vị trí cũ nên độc dịch chuyển d=0.
\item[b)] Tốc độ trung bình của con kiến là: $v_{tb}=\dfrac{s}{t}=\dfrac{18,84}{3}=6,28\ cm/s$.\\
Do $d=0$ nên vận tốc trung bình 0.
\item 			\item 
\end{itemize}
}
\end{vidu}
\begin{vidu}
Một người đi bằng thuyền với tốc độ $2\ m/s$ về phía Đông. Sau khi đi được $2,2\ km$, người này lên ô tô đi về phía Bắc trong $15$ phút với tốc độ $60\ km/h$. 
\begin{listEX}[2]
\item Tính quãng đường đi được.
\item Tính độ dịch chuyển.
\item Tính tốc độ trung bình.
\item Tính vận tốc trung bình.
\end{listEX}
\haidong{13}
%\choice
%{Tổng quãng đường đã đi là $17,2\ km$}
%{Độ dịch chuyển là 15,16 km}
%{Tốc độ trung bình là $8,6\ %m/s$}
%{\True Vận tốc trung bình bằng $8,6\ m/s$}
\loigiai{
15 phút $=0,25\ h$
$2\ m/s=7,2\ km/h$
\begin{itemize}
\item Tổng quãng đường đã đi là $2,2+ 0,25.60=17,2$ km.
\item Độ dịch chuyển là $\sqrt{2,2^2+(0,25.60)^2}=15,16\ km$.
\item Tốc độ trung bình: $v=\dfrac{s}{t}=\dfrac{17,2}{\dfrac{2,2}{7,2}+0,25}=30,96\ km/h=8,6\ m/s.$ 
\item Vận tốc trung bình: $v=\dfrac{d}{t}=\dfrac{15,2}{\dfrac{2,2}{7,2}+0,25}=27,36\ km/h=7,6\ m/s$.
\end{itemize}
}
\end{vidu}




\begin{vidu}%[1P1-2.2K][Loai1]
Một người đi xe máy chuyển động thẳng đều từ $A$ lúc $5$ giờ sáng và tới $B$ lúc $7$ giờ $30$ phút, $AB = 150\ km$. Tốc độ của xe là
\choice
{$40\ km/h$}
{$50\ km/h$}
{\True $60\ km/h$}
{$70\ km/h$}
\haidong{3}
\loigiai{
\begin{itemize}
\item Thời gian đi: \bth{t=7,5-5=2,5h}. 
\item Tốc độ xe lúc đi: 
$v=\dfrac{s}{t}=\dfrac{150}{2,5}=60\ km/h$.
\end{itemize}
}
\end{vidu}

\loai{Tính vận tốc trung bình}
\begin{vidu}
Một người bơi dọc theo chiều dài $100\ m$ của bể bơi hết $60\ s$ rồi bơi quay về lại chỗ xuất phát trong $70\ s$. Trong suốt quãng đường bơi đi và bơi về, vận tốc trung bình của người đó là
\choice
{\True $0\ m/s$}
{$0,769\ m/s$}
{$1,538\ m/s$}
{$1,43\ m/s$}
\haidong{12}
\loigiai{
}
\end{vidu}

\begin{vidu}
\immini[thm]{Trong một lần thử xe ô tô, người ta xác định được độ dịch chuyển của xe tại các thời điểm cách nhau cùng khoảng thời gian $1\ s$ (xem bảng dưới đây). Vận tốc trung bình của ô tô trong $3$ giây cuối cùng là}
{\begin{tabular}{|c|c|c|c|c|c|c|}
\hline 
d (m)&0  &2,3  &9,2  &20,7  &36,8  &57,5  \\ 
\hline 
t (s)&0  &1,0  &2,0  &3,0  &4,0  &5,0  \\ 
\hline 
\end{tabular}} 
\choice
{\True $16,1\ m/s$}
{$2,3\ m/s$}
{$12,27\ m/s$}
{$11,5\ m/s$}
\haidong{12}
\loigiai{
\begin{itemize}
\item Vận tốc trung bình: $v=\dfrac{\Delta d}{\Delta t}$
\item Trong 3 s cuối, ô tô có: \\
Độ dời: $\Delta d=57,5-9,2=48,3\ m$\\
Khoảng thời gian: $\Delta t=3\ s\Rightarrow v=\dfrac{\Delta d}{\Delta t}=\dfrac{48,3}{3}=16,1\ m/s$
\end{itemize}
}
\end{vidu}
\loai{Tốc độ trung bình và vận tốc trung bình}




\loai{Tính tốc độ trung bình theo s và t}
\begin{vidu}%[1P1-2.3B][Loai1]
Một vật chuyển động thẳng không đổi chiều trên quãng đường dài $35\ m$. Nửa quãng đường đầu vật chuyển động thẳng đều hết thời gian $t_1 = 5\ s$, nửa quãng đường sau vật chuyển động thẳng đều hết thời gian $t_2 = 2\ s$. Tốc độ trung bình trên cả quãng đường là
\choice
{$7\ m/s$}
{$6\ m/s$}
{\True $5\ m/s$}
{$4\ m/s$}	
\haidong{12}
\loigiai{
Tốc độ trung bình: $v_{tb}=\dfrac{s}{t}=\dfrac{35}{7}=5\ m/s$.
}
\end{vidu}
\begin{vidu}
Một người bơi dọc theo chiều dài 60 m của bể bơi hết $40\ s$, rồi quay lại về chỗ xuất phát trong $60\ s$. Gọi $v_1, v_2$ và $v_3$ lần lượt là tốc độ trung bình: trong lần bơi đầu tiên theo chiều dài của bể bơi;\ trong lần bơi về và trong suốt quãng đường đi và về. Tổng $\left(v_1+v_2-v_3\right)$ có giá trị là
\choice
{\True $1,3\ m/s$}
{$4,2\ m/s$}
{$3,6\ m/s$}
{$3,5\ m/s$}
\haidong{12}
\loigiai{
\begin{itemize}
\item Lần đi: $v_1=\dfrac{60}{40}=1,5\ m/s$
\item Lần về: $v_2=\dfrac{60}{60}=1\ m/s$
\item Cả đi và về: $v_3=\dfrac{2.60}{40+60}=1,2\ m/s \Rightarrow v_1+v_2-v_3=1,3\ m/s$
\end{itemize}
}
\end{vidu}
\begin{vidu}%[1P1-2.3K][Loai1] 
Một người đi xe đạp đã đi quãng đường $s_1 = 4\ km$ với tốc độ trung bình $v_1 = 12\ km/h$. Sau đó người ấy dừng lại để sửa xe trong $t_2 = 40$  phút rồi đi tiếp quãng đường $s_3 = 8\ km$ với tốc độ trung bình $v_3 = 8\  km/h$. Tốc độ trung bình của người ấy trên tất cả quãng đường đã đi là	
\choice
{$8\ km/h$}
{\True $6\ km/h$}
{$12\ km/h$}
{$10\ km/h$}
\haidong{15}
\loigiai{
\begin{itemize}
\item Thời gian người đó đi trong đoạn đường 4 km đầu: $t_1=\dfrac{s_1}{v_1}=\dfrac{4}{12}=\dfrac{1}{3}\ h$.
\item Khi sửa xe người đó không đi nên $s_2 = 0$  nhưng thời gian sửa xe là: $t_2=40ph=\dfrac{2}{3}\ h$.
\item Thời gian người đó đi trong đoạn đường 8 km sau: $t_3=\dfrac{s_2}{v_2}=\dfrac{8}{8}=1\ h$.
\item Vậy tốc độ trung bình trên toàn bộ quãng đường là:$v=\dfrac{s_1+s_2+s_3}{t_1+t_2+t_3}=6\ km/h$.
\end{itemize}}
\end{vidu}
\loai{Tính tốc độ trung bình theo v và t}
\begin{vidu}
Một người tập thể dục chạy trên một đường thẳng. Lúc đầu người đó chạy với tốc độ trung bình $5\ m/s$ trong thời gian $4$ phút. Sau đó người ấy giảm tốc độ còn $4\ m/s$ trong thời gian $6$ phút. Tốc độ trung bình trong toàn bộ thời gian chạy có giá trị là
\choice
{$3,5\ m/s$}
{$5,6\ m/s$}
{$4,8\ m/s$}
{\True $4,4\ m/s$}
\haidong{12}
\loigiai{
$v_{tb}=\dfrac{s}{t}=\dfrac{v_1 t_1+v_2 t_2}{t_1+t_2}=\dfrac{5 \cdot 4 \cdot 60+4 \cdot 6 \cdot 60}{(4+6) \cdot 60}=4,4\ m/s
$
}
\end{vidu}
\begin{vidu}
Một ô tô chạy trên một đoạn đường thẳng từ địa điểm $A$ đến địa điểm $B$ phải mất một khoảng thời gian $t$. Tốc độ của ô tô trong một phần tư của khoảng thời gian này là $60\ km/h$ và trong phần còn lại là $40\ km/h$. Tốc độ trung bình của ô tô trên cả đoạn đường $AB$ là
\choice
{$48\ km/h$}
{$50\ km/h$}
{$36\ km/h$}
{\True $45\ km/h$}
\haidong{12}
\loigiai{
\begin{itemize}
\item Tốc độ trung bình: $v=\dfrac{s}{t}=\dfrac{s_1+s_2}{t}=\dfrac{v_1 \cdot 0,25 t+v_2 \cdot 0,75 t}{t}=\dfrac{0,25 v_1+0,75 v_2}{1}=45\ km/h$
\end{itemize}
}
\end{vidu}
\loai{Tính tốc độ trung bình theo s và v}
\begin{vidu}%[1P1-2.3K][Loai1]
Một xe đi nửa đoạn đường đầu tiên với tốc độ trung bình $v_1=12\ km/h$ và nửa đoạn đường sau với tốc độ trung bình $v_2=20\ km/h$. Tốc độ trung bình trên cả đoạn đường là
\choice
{$10\ km/h$}
{$12\ km/h$}
{$14\ km/h$}
{\True $15\ km/h$}
\haidong{12}
\loigiai{
\begin{itemize}
\item Thời gian đi nửa đoạn đường đầu: \bth{t_1=\dfrac{s_1}{v_1}=\dfrac{s}{2v_1}=\dfrac{s}{24}}.
\item Thời gian đi nửa đoạn đường sau: \bth{t_2=\dfrac{s_2}{v_2}=\dfrac{s}{2v_2}=\dfrac{s}{40}}.
\item Tốc độ trung bình của xe trên cả đoạn đường: \bth{v=\dfrac{s}{t}=\dfrac{s}{t_1+t_2}=\dfrac{s}{\dfrac{s}{24}+\dfrac{s}{40}}=\dfrac{15s}{s}=15\ km/h.
}
\end{itemize}
}
\end{vidu}
\begin{vidu}
Một người đi xe đạp chuyển động trên một đoạn đường thẳng $AB$ có độ dài là s. Tốc độ của xe đạp trong một phần tư đầu của đoạn đường này là $12\ km/h$, trong một phần năm tiếp theo là $16\ km/h$ và trong phần còn lại là $22\ km/h$. Tốc độ trung bình của xe đạp trên cả đoạn đường $AB$ có giá trị là
\choice
{$48,23\ km/h$}
{$15,15\ km/h$}
{$14,23\ km/h$}
{\True $17,14\ km/h$}
\haidong{12}
\loigiai{
$v_{tb}=\dfrac{s}{t}=\dfrac{s}{t_1+t_2+t_3}=\dfrac{s}{\dfrac{s / 4}{v_1}+\dfrac{s / 5}{v_2}+\dfrac{s\left(1-1/4-1/5\right)}{v_3}}=17,14\ km/h
$
}
\end{vidu}




%mac
\begin{vidu}%[1P1-2.2Y][Loai1]
Một máy bay phản lực chuyển động thẳng đều có tốc độ bằng $2400\ km/h$. Nếu muốn bay liên tục trên khoảng cách $6000\ km$ thì máy bay phải bay trong 
\choice
{$2$ giờ $50$ phút}
{$5$ giờ $20$ phút}
{\True $2$ giờ $30$ phút}
{$3$ giờ $20$ phút}	
\loigiai{
Thời gian để máy bay bay được quãng đường 6000 km là: $t = \dfrac{s}{v} = 2,5\ h = 2h30'$.
}
\haidong{10}
\end{vidu}


\begin{vidu}%[1P1-2.2Y][Loai1]
Một người đi bộ với tốc độ không đổi bằng $1,9\  m/s$. Nếu người đó đi không nghỉ thì sau bao lâu sẽ đến một địa điểm cách nơi xuất phát $780\ m$?
\choice
{$6,50$ phút}
{\True $6$ phút $50$ giây}
{$650$ giây}
{$7$ phút}
\haidong{13}
\loigiai{		
Thời gian để người đó đi hết quãng đường 780 m là: $t=\dfrac{s}{v}=\dfrac{780}{1,9}\approx 410\ s=6\ \text{phút}\ 50\ s$.
}
\end{vidu}

\begin{vidu}%[1P1-2.2K][Loai1]
Một người đi xe máy chuyển động thẳng đều từ $A$ lúc $5$ giờ sáng và tới $B$ lúc $7$ giờ $30$ phút, $AB = 150\ km$. Tới $B$ xe dừng lại $45$ phút rồi đi về $A$ với $v = 50\ km/h$. Xe tới $A$ lúc mấy giờ?
\choice
{$11h45'$}
{$10h45'$}
{\True $11h15'$}
{$9h15'$}
\loigiai{
\begin{itemize}
\item Thời điểm người đó lúc bắt đầu về: \bth{t=7h30'+45'=8h15'}. 
\item Thời gian đi về A: $t=\dfrac{s}{v}=\dfrac{150}{50}=3h$.
\item Vậy xe trở về A lúc \bth{t=8h15'+3h=11h15'}.
\end{itemize}
}
%<MyLT>
\haidong{13}
\end{vidu}
\loai{Tính thời gian ngắn nhất - Câu này đọc rất khó nhưng giải rất dễ}
\begin{vidu}%[1P1-2.2K][Loai1]
\immini[thm]{Một người xuất phát từ $A$ tới bờ sông để lấy nước rồi từ đó mang nước đến $B$. $A$ cách bờ sông một khoảng $AM = 60\ m$; $B$ cách bờ sông một khoảng $BN = 300\ m$. Khúc sông $MN$ dài $480\ m$ và coi là thẳng. Từ $A$ và $B$ tới bất kì điểm nào của bờ sông $MN$ đều có thể đi theo các đường thẳng (hình vẽ). Hỏi muốn quãng đường cần đi là ngắn nhất thì người đó phải đi theo con đường như thế nào và tính chiều dài quãng đường ấy? Nếu người ấy chạy với tốc độ $v = 6\ m/s$ thì thời gian phải chạy hết bao nhiêu?}{\begin{tikzpicture}[draw=Mapcolor]
\tkzInit[ymin=-0.5,ymax=3.5,xmin=-0.5,xmax=5.5]\tkzClip
\tkzDefPoints{0/0/m, 5/0/n, 0/1.5/a, 5/3/b, 5/-3/b', 3.5/0/i, 5/1.5/p}
\tkzInterLL(m,n)(a,b')\tkzGetPoint{j}
\tkzDrawSegments[Mapcolor](m,n)
\fill[pattern=north east lines] (m)rectangle(5,0.1);
\tkzDrawSegments[dashed,Mapcolor](a,m b,n)
\tkzDrawPoints[size=3,fill=white](a,b)
\tkzLabelPoint[above](a){$A$}
\tkzLabelPoint[above](b){$B$}
\tkzLabelPoint[below](m){$M$}
\tkzLabelPoint[below](n){$N$}
\end{tikzpicture}}
\loigiai{
\immini[thm]{\begin{itemize}
\item Giả sử người đi theo đường AIB. Gọi B' là điểm đối xứng của B qua bãi sông.
\item Ta có: $AIB=AI+IB=AI+IB'=AIB'$
\item Để AIB ngắn nhất thì ba điểm A, I, B' thẳng hàng. Lúc đó I trùng J.
\item Từ hình vẽ ta có:
$\begin{cases}
AP=MN=480\ m\\ 
B'P=B'N+NP=360\ m\\ 
\end{cases}$ $\Rightarrow AB'=\sqrt{AP^2+B'P^2}=600\ m$
\item Thời gian ngắn nhất: $t_{min}=\dfrac{AB'}{v}=\dfrac{600}{6}=100\ s$.
\end{itemize}}{\begin{tikzpicture}
\tkzInit[ymin=-3.5,ymax=3.5,xmin=-0.5,xmax=5.5]\tkzClip
\tkzDefPoints{0/0/m, 5/0/n, 0/1.5/a, 5/3/b, 5/-3/b', 3.5/0/i, 5/1.5/p}
\tkzInterLL(m,n)(a,b')\tkzGetPoint{j}
\tkzDrawSegments[](m,n)
\fill[pattern=north east lines] (m)rectangle(5,0.1);
\tkzDrawSegments[dashed](a,m b,n a,i i,b a,p b,b' j,b' i,b')
\draw[very thick,postaction={decorate, decoration={markings,
mark=at position 0.5 with {\arrow{>};}
}}]
(a)--(j);
\draw[very thick,postaction={decorate, decoration={markings,
mark=at position 0.5 with {\arrow{>};}
}}]
(j)--(b);
\tkzDrawPoints[size=3,fill=white](a,b,i,j,b',p)
\tkzLabelPoint[above](a){$A$}
\tkzLabelPoint[right](b){$B$}
\tkzLabelPoint[right](b'){$B'$}
\tkzLabelPoint[right](p){$P$}
\tkzLabelPoint[below](m){$M$}
\tkzLabelPoint[below right](n){$N$}
\tkzLabelPoint[below left](i){$I$}
\tkzLabelPoint[below left](j){$J$}
\end{tikzpicture}}
}
\haidong{13}
\end{vidu}


